% Part: proof-theory
% Chapter: normalization
% Section: normalization-thm

\documentclass[../../../include/open-logic-section]{subfiles}

\begin{document}

\olfileid{pt}{nor}{thm}

\olsection{正规化定理}

如果一个推导没有切割段，则它处于\emph{正规}形式。显然，当且仅当既不能对一个推导应用置换转换，也不能应用约减转换时，它才是正规的。对一个推导进行\emph{正规化}，是指通过反复应用置换转换和约减转换，直至不再有切割存在，从而将其转化为一个正规证明。这种转化总是可能的，这就是\emph{正规化定理}。

\begin{thm}\ollabel{thm:normalization} 任何 \Log{N1i} 推导 $\delta$ 都是可正规化的，即，它可以被转化为一个没有切割的推导~$\delta'$。
\end{thm}

\begin{proof}
  对 $\delta$ 的切割秩和切割长度施归纳。归纳假设为：任何切割秩更低，或切割秩相同但切割长度更小的推导，都是可正规化的。

  若切割长度为~$0$，则不存在切割，$\delta$ 已经是正规的。因此，设 $\delta$ 的切割秩和切割长度均~$> 0$。选取一个最上方、最右的极大切割段。若该切割段长度 $>1$，则应用置换转换。若该切割段长度为~$1$，则应用相应的约减转换。由此将得到一个新的推导，它要么切割秩相同但切割长度更小，要么切割秩更低。于是归纳假设适用。
\end{proof}

上述证明中使用的策略（总是处理最上方、最右的极大切割段）提示了一种具体的转换顺序，必须按此顺序对推导应用转换才能得到正规推导。但令人惊讶的是，顺序其实无关紧要。以任意顺序应用置换转换和约减转换，最终都能得到一个正规证明。

段和切割的定义很容易为 \Log{N2i} 给出表述，而且适用于 \Log{N1i} 的置换转换和约减转换可以直接转化为适用于 \Log{N2i} 的转换。因此，正规化定理对 \Log{N2i} 同样成立。

\begin{thm}\ollabel{thm:normalization-N2i} 任何 \Log{N2i} 推导 $\delta$ 都是可正规化的，即，它可以被转化为一个没有切割的推导~$\delta'$。
\end{thm}

\end{document}